% =============================================================================
\section{Conclusiones}
\label{sec:conclusiones}
% =============================================================================

Se ha presentado un estudio sistemático mediante dinámica molecular del rebote
inelástico de clústeres Lennard-Jones sobre una superficie cristalina FCC de
argón, cubriendo cinco tamaños de clúster, cinco velocidades de impacto, tres
temperaturas superficiales, dos energías de adhesión y tres ángulos de incidencia
en \num{1350} simulaciones independientes. Las principales conclusiones son las
siguientes.

\medskip\noindent
\textbf{(i) Partición tridimensional de la energía.}
La energía cinética normal incidente $\Einc$ se distribuye entre tres canales
cuantificables: rebote ($\cor^2\,\Einc$), modos internos del clúster
($\alphaint\,\Einc$) y fonones de la red ($\alphasurf\,\Einc$), con balance
verificado numéricamente con precisión de doble precisión. Para un proyectil
monoatómico ($N=1$), toda la disipación inelástica va a la red
($\alphaint \equiv 0$), con coeficientes de acomodación energética
$\alphasurf \in \TODO{[0{,}15,\, 0{,}75]}$ consistentes con la literatura.

\medskip\noindent
\textbf{(ii) El canal interno crece con el tamaño y domina para clústeres grandes.}
$\alphaint$ aumenta desde cero para $N=1$ hasta $\alphaint \approx \TODO{0{,}55}$
para $N=55$ a velocidades altas, siguiendo un crecimiento rápido para $N \leq 7$
y una saturación progresiva para $N \geq 13$. Este comportamiento es consistente
con el aumento de la densidad de modos vibracionales de baja frecuencia del
clúster que pueden acoplarse al impulso de contacto. Para $N \geq 13$, el canal
interno supera al canal superficial como mecanismo dominante de disipación.

\medskip\noindent
\textbf{(iii) Ley de potencias con exponente dependiente del tamaño.}
El COR sigue $\cor \sim A\,v_0^{-\beta}$ con exponentes que crecen de
$\beta_{N=1} \approx \TODO{0{,}30}$ a $\beta_{N=55} \approx \TODO{0{,}55}$.
Todos los valores superan la predicción $\beta=1/5$ del modelo continuo
viscoelástico (Kuwabara-Kono), atribuible a la anharmonicidad del potencial
LJ y a la naturaleza discreta de la red. La ausencia de anomalías en los
números mágicos icosaédricos ($N=13,\, 55$) indica que la estabilidad
estructural extra de estos clústeres no modifica el flujo de energía durante
el impacto en el régimen de velocidades intermedias estudiado.

\medskip\noindent
\textbf{(iv) El número de Weber colapsa la transición rebote-adhesión.}
Las curvas de fracción de sticking para distintos $(N, \epsps)$ colapsan
sobre una curva maestra en función de $\We = \Einc/(N\,\epsps)$, con umbral
$\We_c \approx \TODO{1{,}5}$ consistente con resultados previos para clústeres
LJ~\cite{awasthi2007}. La temperatura superficial introduce una dispersión
residual no contenida en $\We$, sugiriendo que una versión extendida del número
de Weber que incluya el efecto térmico mejoraría el colapso.

\medskip\noindent
\textbf{(v) Efectos de $\epsps$ y $\Tstar$.}
Reducir la energía de adhesión proyectil-superficie de $\epsps = 1{,}00\,\varepsilon$
a $0{,}30\,\varepsilon$ eleva el COR principalmente a bajas velocidades, mientras
que el canal interno $\alphaint$ permanece prácticamente inalterado, lo que
confirma que éste refleja propiedades intrínsecas del clúster. Aumentar $\Tstar$
reduce el COR en toda condición y amplía el canal interno para clústeres con
$N \geq 7$, en acuerdo con la predicción del modelo GLE de acoplamiento
proyectil-fonones~\cite{tully1980}.

\medskip
Los resultados establecen un cuadro cuantitativo del flujo de energía durante
el impacto de nanoclústeres LJ sobre FCC, con implicaciones para el diseño de
procesos de deposición de nanopartículas donde el control del régimen
rebote-adhesión y del calentamiento del clúster es relevante~\cite{haberland1995,
cheng1993}.

\subsection*{Trabajo futuro}

Las extensiones naturales de este trabajo incluyen:

\begin{itemize}

    \item \textbf{Separación del canal rotacional.}
    Calcular el tensor de inercia frame a frame para descomponer
    $\alphaint = \alpha_\mathrm{vib} + \alpha_\mathrm{rot}$
    en impactos oblicuos, completando el balance energético vectorial.

    \item \textbf{Potenciales realistas.}
    Repetir el barrido con potenciales EAM para metales (Cu, Au) o AIREBO
    para carbono, donde la comparación con experimentos de deposición de
    nanoclústeres en masa seleccionada es directa~\cite{haberland1995}.

    \item \textbf{Clústeres más grandes.}
    Extender el estudio a $N \in \{147,\, 309\}$ (tercer y cuarto icosaedros
    de Mackay) para determinar si la saturación de $\alphaint$ continúa o
    emerge un nuevo régimen en el que el clúster se comporta como un
    proyectil sólido macroscópico.

    \item \textbf{Número de Weber extendido.}
    Determinar el coeficiente $\alpha$ en $\We^* = \Einc / [N(\epsps + \alpha k_BT)]$
    que minimice la dispersión de la curva maestra al variar $\Tstar$.

    \item \textbf{Efectos cuánticos a bajas temperaturas.}
    Para $\Tstar < 0{,}05$ en argón ($T < 5$~K), las correcciones de
    cero-punto a los fonones de la red pueden modificar cuantitativamente
    el canal $\alphasurf$. Su estimación mediante DM con corrección de camino
    de integrales (PIMD) es factible con el mismo sistema modelo.

\end{itemize}
